\subsection{Conditional arithmetic gluing of repairs in increasing form degree}

\begin{theorem}[C30: \status{Proved under stated hypotheses}; a common-source reduction]
\label{sa08:C30}
Let $k$ be a commutative ring, put $S=\Gamma_k(T)$, and let
$A,B$ be admissible generalized theories. Let $P\to A,B$ be
the common partial model of Statement~\ref{sa08:C21}.
For each prime $p$, write $k_p=k\otimes_\Z\Z_{(p)}$ and
$P_p=P\otimes_k k_p$, and similarly for the targets.

Suppose that for every prime $p$ there are an admissible
theory $E_p$ over $k_p$ and compatible maps
\[
 P_p\longrightarrow E_p\longrightarrow A_p,B_p.
\]
Suppose the first map is an isomorphism in all form degrees
$0\leq q<d_p$, in every weight and simplex, where the integers
$d_p\geq0$ tend to infinity with $p$: for each $d$, only
finitely many primes satisfy $d_p\leq d$.

Then there is an admissible common source $H\to A,B$ over
$k$. The condition $q\leq w$, if imposed on all the data,
is preserved. No termwise flatness hypothesis is needed
for this assertion. If $k=\Z$, the targets and all the
local repairs have flat terms, then $H$ can also be taken
termwise flat over $\Z$.

This is a conditional common-source and connectivity
theorem. It does not construct the assumed repairs and
does not assert that every loop in the localized category
is trivial.
\end{theorem}
\begin{proof}
\par\medskip\noindent\textbf{A pure arithmetic square.}
The classical completion version is Sullivan's arithmetic square
\cite[Chapter~1, Propositions~1.17--1.18, pp.~28--30]{A_sullivan_mit2005}.
We prove here the localization-product variant needed for arbitrary
modules and for the uniform form-degree construction below.

Set
\[
 R=\prod_{p\ \mathrm{prime}}\Z_{(p)},
 \qquad R_\Q=R\otimes_\Z\Q.
\]
We first prove the short exact sequence
\begin{equation}
 0\longrightarrow\Z\longrightarrow\Q\oplus R
   \xrightarrow{\ (a,b)\mapsto b-a\ }R_\Q
   \longrightarrow0.
 \tag{D-AR1}
\end{equation}
The kernel is $\Q\cap\prod_p\Z_{(p)}=\Z$.
An element of $R_\Q$ has the form $b/D$ for one positive
integer $D$ and one $b=(b_p)\in R$. Only the finitely many
primes dividing $D$ can give nonintegral local components.
Choose $c\in\Z$ with
\[
 c\equiv b_p\pmod{p^{v_p(D)}\Z_{(p)}}
 \qquad(p\mid D),
\]
using the finite Chinese remainder theorem. Then
$b/D-c/D\in R$, and the pair
$(-c/D,b/D-c/D)$ maps to $b/D$, proving surjectivity. The quotient
$R_\Q$ is flat over $\Z$, so (D-AR1) is pure. Tensoring
with any abelian group is therefore exact.

Put $k_R=k\otimes_\Z R$, $k_\Q=k\otimes_\Z\Q$, and
$k_{R\Q}=k\otimes_\Z R_\Q$. For every $k$-module $M$,
purity gives
\begin{equation}
 0\longrightarrow M\longrightarrow M_\Q\oplus M_R
       \longrightarrow M_{R\Q}\longrightarrow0.
 \tag{D-AR2}
\end{equation}
All these scalar extensions are flat. Here $M_R$ means
$M\otimes_\Z R$; it is not replaced by
$\prod_p M_p$. This distinction is essential for arbitrary
modules and for the uniformity below.

\par\medskip\noindent\textbf{An actual repair over the product coefficient ring.}
Let $e_p\in R$ be the coordinate idempotent, so
$e_pk_R=k_p$. Choose increasing finite sets
\[
 F_q=\{p:d_p\leq q\}\ \cup\ \{p:p\leq q+2\},
 \qquad e_{F_q}=\sum_{p\in F_q}e_p.
\]
Define the term of a bigraded simplicial $k_R$-module $E_R$
by
\begin{equation}
 E_R^{w,q}
  =(1-e_{F_q})P_R^{w,q}
       \oplus\bigoplus_{p\in F_q}E_p^{w,q}.
 \tag{D-AR3}
\end{equation}
This is a finite idempotent decomposition in each form degree.
It is not the unrestricted product of the local cochain
modules. If a larger finite set $F\supset F_q$ is used in
(D-AR3), every new component has $q<d_p$, so the given map
$P_p^{w,q}\to E_p^{w,q}$ is an isomorphism. Thus refinement
gives a canonical identification of the same module.

Define the differential after refining the source and target
to $F_{q+1}$. Define a product of form degrees $q,q'$ after
refining to $F_{q+q'}$, and define an $n$th divided power
after refining to $F_{nq}$. On each selected component use
the corresponding operation of $E_p$; on the complementary
idempotent use the operation of $P_R$. All identities hold
after a common finite refinement, where they are the
identities of those algebras. This defines a strict PD
differential $S_{k_R}$-algebra, including its simplicial
operators, coefficient map, and augmentation. The maps
\[
 P_R\longrightarrow E_R\longrightarrow A_R,B_R
\]
are defined in the same way and are compatible.

We verify admissibility with uniform bounds. In a fixed form
range $0\leq q\leq d$, refine all terms to $F_d$.
The complementary summand is a flat scalar extension of
$P$, whose rows satisfy eventual horizontal acyclicity by
Statement~\ref{sa08:C21}. There are only finitely many
other summands, and they are rows of the admissible $E_p$.
The maximum of their finitely many thresholds gives the
horizontal bound for $E_R$.

For vertical cohomology through degree $d$, refine through
degree $d+1$ to $F_{d+1}$. The complementary coefficient
ring inverts every prime at most $d+3$, hence inverts
$\lfloor(d+1)/2\rfloor!$. Statement~\ref{sa08:C21}
therefore makes the complementary $P$ summand a coefficient
resolution through degree $d$ in every simplex and weight.
For the finitely many selected primes use the vertical
eventual bounds of $E_p$, and take their maximum. This
proves the required vertical assertion, uniformly in the
simplex. Thus $E_R$ is admissible. Also $P_\Q$ is
admissible by the same result about $P$.

\par\medskip\noindent\textbf{Gluing the two coefficient regions.}
Write $D=(E_R)_\Q$. There is a compatible map
$r\colon P_\Q\to D$ induced by $P_R\to E_R$.
Use the simplicial interval $I=\Delta[1]$ and form
\begin{equation}
 E=P_\Q\mathop{\times}_{D,\operatorname{ev}_0}D^I
       \mathop{\times}_{\operatorname{ev}_1,D}E_R.
 \tag{D-AR4}
\end{equation}
The cotensor and relative-path facts of
Statement~\ref{sa08:C16} apply over $k_{R\Q}$.
In each fixed row, all three theories become horizontally
acyclic in sufficiently large weight. Endpoint restriction
$D^I\to D\oplus D$ is then surjective, and its
both-endpoint kernel is horizontally acyclic. Hence (D-AR4)
is horizontally acyclic in that range; a finite maximum
works for any fixed form range.

For the vertical assertion, choose weights large enough
for the finite-range comparisons of those three theories
and their interval cotensors. In each simplex the kernel
of both endpoints has only coefficient cohomology
$S_{k_{R\Q}}^w$ in degree one. The endpoint theories have
coefficient cohomology $S_{k_\Q}^w\oplus S_{k_R}^w$
in degree zero. The connecting map is the difference map
\[
 S_{k_\Q}^w\oplus S_{k_R}^w
       \longrightarrow S_{k_{R\Q}}^w.
\]
By (D-AR2), its kernel is $S_k^w$ and its cokernel is
zero. The long exact sequence proves that $E$ resolves
$S_k$ through the required form range. The path comparison
is uniform in the simplex by Statement~\ref{sa08:C16}.
Thus $E$ is admissible. Its augmentation is also defined
over $k$: after augmentation the path is constant, and
the two endpoint values descend uniquely by (D-AR2).

\par\medskip\noindent\textbf{Obtaining actual maps to the original targets.}
For an admissible theory $C$ over $k$, define
\[
 J(C)=C_\Q\mathop{\times}_{C_{R\Q},\operatorname{ev}_0}
       (C_{R\Q})^I
       \mathop{\times}_{\operatorname{ev}_1,C_{R\Q}}C_R.
\]
There is a natural coefficient-preserving map
$\eta_C\colon C\to J(C)$ using constant paths.
In every weight and form degree one has a short exact
sequence of simplicial modules
\[
 0\longrightarrow C\xrightarrow{\eta_C}J(C)
       \longrightarrow P_0(C_{R\Q})\longrightarrow0,
\]
where $P_0$ denotes paths vanishing at the initial endpoint.
The last map subtracts the constant initial value from the
path. Its surjectivity follows from (D-AR2), applied to
the final value of any based path. The kernel is exactly
$C$ by that same arithmetic sequence.

The based-path module is horizontally contractible in every
row, and its vertical complex is eventually acyclic by
Statement~\ref{sa08:C16}. Consequently $J(C)$ is admissible,
and $\eta_C$ is a horizontal weak equivalence in every
weight and form degree. The compatibility of
$P_p\to E_p\to A_p,B_p$ gives actual maps
\[
 E\longrightarrow J(A),\qquad E\longrightarrow J(B).
\]
No choices of rational homotopies are needed here: the
rational endpoint diagrams commute exactly.

For clarity, we replace this zigzag by a genuine common
source. Factor $\eta_C$ by the simplicial mapping path
\[
 \widetilde C=C\mathop{\times}_{J(C),\operatorname{ev}_0}
                     J(C)^I
       \xrightarrow{\operatorname{ev}_1}J(C).
\]
Since $\eta_C$ is a horizontal weak equivalence, the last
map is a trivial fibration in every simplicial module row.
It is degreewise surjective and has horizontally acyclic
kernel. The projection $\widetilde C\to C$ has a based-path
kernel, so $\widetilde C$ is admissible. Its kernel over
$J(C)$ is also vertically eventually acyclic, since the
surjection preserves the coefficient resolution.
Now put
\[
 H=E\mathop{\times}_{J(A)\times J(B)}
              (\widetilde A\times\widetilde B).
\]
The projection $H\to E$ is degreewise surjective; its
kernel is the direct sum of the two just described
kernels. It is horizontally acyclic in every row and
vertically eventually acyclic. Hence $H$ is admissible
and its projections give the required maps to $A$ and $B$.

All constructions preserve both gradings; interval paths
are simplicial cotensors, so no new form-degree variable
is introduced. This proves the optional support assertion.
If $k=\Z$ and the terms in the statement are flat, the
finite idempotent construction is flat over $R$. Every
rational path term is a $\Q$-vector space. The subsequent
fiber products and kernels are subgroups of torsionfree
abelian groups, and are therefore torsionfree and flat
over $\Z$. This proves the stated flat version. The same
argument is not used to assert flatness over an arbitrary
coefficient ring.
\end{proof}

\par\medskip\noindent\textbf{The remaining hypotheses are substantive.}
The required $P_p\to E_p$ must retain the low form degrees
as actual isomorphisms, uniformly in all weights and
simplices, and its maps to both targets must extend those
of $P$. An isolated acyclic disk closure does not provide
these relative maps. Statement~\ref{sa08:C30} removes the
additional infinite-prime gluing step once such repairs
exist; it does not prove their existence. Nor do the
assumed maps supply coherent homotopies for parallel
arrows. The terminal-localization assertion therefore
remains separate from this common-source reduction.
